\documentclass[11pt]{article}
\usepackage{climbingguide}

% Local image directory mapping bypassed for isolated testing
% \graphicspath{{C:/Users/lucas/Downloads/extracted_images/}}

\begin{document}

% --- PAGE 1: COVER ---
\CGCoverPage{example-image.jpg}{IBICOARA}{22º EENe | 2025}

% --- PAGE 2: APRESENTAÇÃO ---
\CGTextSingleColumn{Expediente}{
    Guia de Escalada de Ibicoara, Bahia - 22º EENe | 2025\\
    Foto da Capa: Ygor Tinoco\\
    Textos: Janine Falcão, Allan Carlos
}

\CGDoubleTextSection{Apresentação}{
    Este é mais um guia para o Encontro de Escaladores do Nordeste! Ele foi elaborado para ajudar e facilitar o acesso dos escaladores aos diversos setores de escalada, que, desta vez, estarão na comunidade de Campo Redondo, em Ibicoara.
}{}{
    Este EENe tem um valor único na verdade, todas as edições são especiais! Mas esta, em particular, além de marcar um destino inédito para o nosso querido encontro, não poderia ser diferente: sendo na Bahia, tinha que ser longe, né?
}
\clearpage

% --- PAGE 3: ZONA BOQUEIRÃO ---
\CGZoneHeader{example-image-a.jpg}{ZONA BOQUEIRÃO}

\CGDoubleTextSection{Setores}{
    \textbf{Mocozado, Mocozinho e Mocó.}\\[4pt]
    A Zona Boqueirão é composta por três setores de escalada distintos: Mocó, Mocozinho e Mocozado. Cada um apresenta características próprias, somando mais de 40 vias, com graus de dificuldade que vão do 5º ao 8º.\\[4pt]
    O acesso a este setor é realizado através da propriedade do Sr. Adjacy. A caminhada leva aproximadamente 20 minutos e conecta os três setores.
}{Betas da zona}{
    O horário mais recomendado para escalar nesta zona é à tarde, com a possibilidade de iniciar a escalada a partir das 11 horas.\\[4pt]
    \CGObsBox{\textbf{\color{brand_color}OBS.} Em algumas épocas do ano, é comum encontrar, na Chapada Diamantina, uma espécie de vespa que se abriga nas fendas das vias. A boa notícia é que essas vespas não apresentam comportamento agressivo.}
}
\clearpage

% --- PAGE 4: MOCOZADO ---
\CGSectorHeader{guide_red}{MOCOZADO}

\CGTextSingleColumn{}{
    O setor Mocozado conta com apenas 2 vias, sendo as mais desafiadoras da Zona, ideais para quem busca maiores dificuldades. O acesso a esse setor é feito por uma pequena chaminé.
}

\CGTopoImage{example-image-b.jpg}

\begin{multicols}{2}
    \CGRoute{GUERRA E PAZ}{Proj.}{30m}{10+2}{Messias e Paulinho}
    \CGRoute{ESCAFÓIDE}{8a}{30m}{13+2}{Messias e Paulinho}
\end{multicols}
\clearpage

% --- PAGE 5: MOCOZINHO ---
\CGSectorHeader{guide_yellow}{MOCOZINHO}

\CGDoubleTextSection{}{
    Diferente do Mocó, o setor Mocozinho é mais acessível, com vias que variam entre 5º e 6sup, todas em proteções fixas e bem protegidas, com destaque para a via "Chama os Amigos".
}{}{
    A base é confortável, mas é necessário tomar cuidado com algumas áreas mais expostas. Uma vantagem desse setor é que é possível escalar até em dias chuvosos.
}

\CGTopoImage{example-image-c.jpg}

\begin{multicols}{2}
    \CGRoute{CAFÉ SEM TABACO}{6sup}{15m}{5+2}{Allan Carlos e Vinicin}
    \CGRoute{EQUILÍBRIO MALUCO}{7c}{15m}{4+2}{Allan Carlos e Vinicin}
    \CGRoute{A BUSCA CONTINUA}{7a}{15m}{5+2}{Allan Carlos e Vinicin}
    \CGRoute{NÃO SUFOQUE O ARTISTA}{5sup}{15m}{4+2}{Allan Carlos e Vinicin}
    \CGRoute{AÇÃO E REAÇÃO}{5sup}{15m}{4+2}{Allan Carlos e Vinicin}
    \CGRoute{DÍZIMO DA PEDRA}{6sup}{17m}{5+2}{Allan Carlos e Vinicin}
    \columnbreak
    \CGRoute{CONVITE ESPIRITUAL}{6sup}{15m}{6+2}{Allan Carlos e Vinicin}
    \CGRoute{GLAUCOMA}{7a}{15m}{6+2}{Allan Carlos e Vinicin}
    \CGRoute{CHAMA OS AMIGOS}{6sup}{16m}{6+2}{Allan Carlos e Vinicin}
    \CGRoute{SECOS E MOLHADOS}{6sup}{16m}{6+2}{Allan Carlos e João Paulo}
    \CGRoute{CARANGA}{6sup}{15m}{7+2}{Allan Carlos e Vinicin}
\end{multicols}

\end{document}
